% 一般规范场路径积分、Faddeev--Popov 规范固定与 BRST 结构
局域规范对称性的出发点可以先用一个简单事实理解：同一个物理内部状态，可以在每个时空点选择不同的内部基底来描述。若这个基底变化依赖于位置，普通导数会额外微分到基底本身，于是两个相邻点上的内部向量不能直接比较。为补偿这种“局域基底变化”而引入的规范势称为\emph{规范联络}；协变导数就是利用它规定相邻时空点之间怎样比较内部状态。电磁势 $A_\mu$ 对应阿贝尔 $U(1)$ 联络；当内部变换不再彼此交换时，同一局域比较原则给出非阿贝尔联络。

量子理论还遇到第二个问题：不同的联络函数可以只相差一次局域规范变换，却表示完全相同的物理状态。把“所有可能的场函数”整体看成一个巨大空间，称为\emph{场空间}；其中一个“点”不是普通坐标，而是一整套函数 $\mathscr C_\mu^a(x)$。从某个场构形出发施加所有可能的局域规范变换，会在场空间中扫出一族彼此物理等价的点，这一族称为\emph{规范轨道}。若路径积分把轨道上的每一点都积分一次，就会把同一个物理状态无限重复计数。

因此规范固定不是增加新的动力学，而是在每条规范轨道上选择一个代表。几何上，这相当于在场空间中选一张横截轨道的“切片”。换变量时还必须乘上 Jacobian；这里的\emph{泛函 Jacobian}只是普通多变量积分 Jacobian 行列式在无限多个场自由度上的对应物。它衡量规范参数发生微小变化时，规范条件函数改变多少。这个 Jacobian 可写成明确的泛函行列式；其标准名称是 Faddeev--Popov 行列式。

令紧致李群 $G$ 表示内部规范变换群，生成元与结构常数沿用第一章的 Lie 代数定义。所有联络构成的场空间记为 $\mathscr A$，一条规范轨道记为 $[\mathscr C]=\{\mathscr C^U\}$。作用量沿轨道不变，因此二次变分在轨道切向上有零模；未经规范固定时，传播子需要反演的二次核便不可逆。采用几何联络归一化
\begin{equation}
 \mathscr C_\mu=\mathscr C_\mu^aT^a,
 \qquad
 D_\mu=\partial_\mu+\ii g_{\rm G}\mathscr C_\mu,
 \label{eq:generic-gauge-cov-derivative-dp11}
\end{equation}
其中 $g_{\rm G}$ 是与一般规范群 $G$ 对应的无量纲规范耦合，$[\mathscr C_\mu]=\mathrm{m^{-1}}$。场强由协变导数的对易子定义，
\begin{align}
 [D_\mu,D_\nu]&=\ii g_{\rm G}\mathcal F_{\mu\nu},\\
 \mathcal F_{\mu\nu}^a
 &=\partial_\mu\mathscr C_\nu^a-\partial_\nu\mathscr C_\mu^a
 -g_{\rm G} f^{abc}\mathscr C_\mu^b\mathscr C_\nu^c,
 \label{eq:generic-field-strength-dp11}
\end{align}
对应的纯规范场拉格朗日密度为
\begin{equation}
 \Lag_{\rm YM}
 =-\frac{\hbar c}{4}\mathcal F_{\mu\nu}^a\mathcal F^{a\mu\nu}.
 \label{eq:generic-ym-lagrangian-dp11}
\end{equation}

一般局域规范变换首先作用在联络本身。与上面的协变导数约定相容的有限变换为
\begin{equation}
\boxed{
 \mathscr C_\mu'
 =U\mathscr C_\mu U^{-1}
 +\frac{\ii}{g_{\rm G}}(\partial_\mu U)U^{-1}.}
\label{eq:generic-finite-gauge-transform-dp68}
\end{equation}
现在才取无穷小形式 $U=I+\ii\vartheta^aT^a+\order(\vartheta^2)$。它给出规范轨道在场空间中的切向量
\begin{equation}
 \boxed{
 \delta\mathscr C_\mu^a
 =-\frac1{g_{\rm G}}(\mathcal D_\mu^{\rm adj})^{ab}\vartheta^b,
 \qquad
 (\mathcal D_\mu^{\rm adj})^{ab}
 =\delta^{ab}\partial_\mu-g_{\rm G}f^{acb}\mathscr C_\mu^c.}
 \label{eq:generic-infinitesimal-gauge-dp11}
\end{equation}
同一规范轨道上的场构形给出相同物理状态，因此沿\eqnrefs{eq:generic-infinitesimal-gauge-dp11} 的切向方向积分只是在重复计算同一个物理点。

选取规范条件 $F^a[\mathscr C]=\omega^a$。它的作用只是从同一条规范轨道上选取一个代表。要把这一步放进路径积分，需要知道规范参数变化会把 $F$ 改变多少。把规范条件对无穷小规范参数的函数导数看成一个线性算符，它就是规范轨道坐标到规范条件坐标之间的 Jacobian 核；其泛函行列式给出积分测度的修正：
\begin{equation}
\boxed{
 \mathcal M_{\rm FP}^{ab}(x,y)
 =\left.\frac{\delta F^a[\mathscr C^\vartheta](x)}{\delta\vartheta^b(y)}\right|_{\vartheta=0},
 \qquad
 \Delta_{\rm FP}[\mathscr C]=\det\mathcal M_{\rm FP}.}
\label{eq:fp-jacobian-definition-dp68}
\end{equation}
算符 $\mathcal M_{\rm FP}$ 称为\emph{Faddeev--Popov 算符}，$\Delta_{\rm FP}$ 称为 Faddeev--Popov 行列式。若这个规范条件在所讨论的局部场空间区域内与每条规范轨道横截一次，函数空间的换元公式便给出
\begin{equation}
 \boxed{
 1=\Delta_{\rm FP}[\mathscr C]
 \int\mathcal D\vartheta\,
 \delta\!\left[F(\mathscr C^\vartheta)-\omega\right].}
 \label{eq:fp-identity-general-dp11}
\end{equation}
这里 $\Delta_{\rm FP}$ 是从规范参数 $\vartheta$ 到规范条件 $F$ 的泛函雅可比，它只修正规范轨道上的积分测度。对协变规范 $F^a=\partial^\mu\mathscr C_\mu^a$，相应算符为
\begin{equation*}
 \mathcal M_{\rm FP}^{ab}
 =-\partial^\mu(\mathcal D_\mu^{\rm adj})^{ab}
\end{equation*}
（略去与场无关的整体常数）。对辅助函数 $\omega^a$ 作高斯平均后得到
\begin{equation}
 \boxed{
 \Lag_{\rm gf}
 =-\frac{\hbar c}{2\xi}
 (\partial^\mu\mathscr C_\mu^a)^2,}
 \label{eq:general-covariant-gf-dp11}
\end{equation}
其中 $\xi$ 是无量纲规范参数。

Faddeev--Popov 行列式仍是一个非局域的泛函行列式。前面已经见过 Grassmann 高斯积分会产生普通矩阵的行列式；把同一个代数恒等式应用到这里，可以用两组 Grassmann 奇的标量场 $c^a,\bar c^a$ 把该行列式写成局域作用量。它们称为\emph{鬼场}（ghost fields）。鬼场是表示规范轨道 Jacobian 的辅助 Grassmann 变量；BRST 物理态条件把它们排除在可探测的渐近态之外。由此得到鬼场拉格朗日量
\begin{equation}
 \boxed{
 \Lag_{\rm gh}
 =\hbar c\,\bar c^a
 \left[-\partial^\mu(\mathcal D_\mu^{\rm adj})^{ab}\right]c^b.}
 \label{eq:general-ghost-lagrangian-dp11}
\end{equation}
阿贝尔理论中 $f^{abc}=0$，$\mathcal M_{\rm FP}$ 与规范场无关，鬼场行列式可并入归一化；非阿贝尔理论中 $\mathcal M_{\rm FP}$ 含 $\mathscr C_\mu$，鬼场因而与规范场相互作用。鬼场不是物理渐近粒子，它们是规范轨道雅可比的局域场表示。

\begin{exbox}[自由 $U(1)$ 场中的规范零模]
一般规范轨道结构在自由电磁场中退化为线性零模。把 Fourier 波矢记为 $\kappa^\mu$，Maxwell 二次核为
\begin{equation}
 \boxed{
 \mathcal K_0^{\mu\nu}(\kappa)
 =-\kappa^2\eta^{\mu\nu}+\kappa^\mu\kappa^\nu,
 \qquad
 \mathcal K_0^{\mu\nu}\kappa_\nu=0.}
 \label{eq:u1-gauge-zero-mode-dp50}
\end{equation}
因此任意 $\mathscr C_\mu=\lambda\kappa_\mu$（$\lambda$ 为任意标量 Fourier 系数）都满足 $\mathcal K_0^{\mu\nu}\mathscr C_\nu=0$，它们正是规范轨道切向的零模。加入协变规范固定项 $-(\hbar c/2\xi)(\partial\!\cdot\!\mathscr C)^2$ 后，核变为
\begin{equation}
 \boxed{
 \mathcal K_\xi^{\mu\nu}
 =-\kappa^2\eta^{\mu\nu}
 +\left(1-\frac1\xi\right)\kappa^\mu\kappa^\nu,}
 \label{eq:u1-gauge-fixed-kernel-dp50}
\end{equation}
其横向与纵向本征值分别为 $-\kappa^2$ 与 $-\kappa^2/\xi$。这个线性阿贝尔例子展示一般规范轨道切片的最简单实现；Faddeev--Popov Jacobian 则由一般规范变换在规范条件上的泛函导数定义。


\medskip

\eqnrefs{eq:u1-gauge-zero-mode-dp50} 已经把问题具体化为“二次核沿 $\kappa^\mu$ 方向有零本征值”。这里加入协变规范固定项，并在横向、纵向投影上重新求二次核的本征值，从而得到可逆的\eqnrefs{eq:u1-gauge-fixed-kernel-dp50}。这一计算只涉及自由 $U(1)$ 场，还不需要鬼场或 BRST。

取 Fourier 约定 $\partial_\mu\to-\ii\kappa_\mu$。自由场强为
\begin{equation*}
 \mathcal F_{\mu\nu}
 =-\ii(\kappa_\mu\mathscr C_\nu-\kappa_\nu\mathscr C_\mu).
\end{equation*}
把它代入 $-\hbar c\,\mathcal F_{\mu\nu}\mathcal F^{\mu\nu}/4$，并把二次型写成
\begin{equation*}
 \frac{\hbar c}{2}\,
 \mathscr C_\mu(-\kappa)\mathcal K_0^{\mu\nu}(\kappa)\mathscr C_\nu(\kappa),
\end{equation*}
得到
\begin{equation*}
 \mathcal K_0^{\mu\nu}
 =-\kappa^2\eta^{\mu\nu}+\kappa^\mu\kappa^\nu.
\end{equation*}
右乘 $\kappa_\nu$ 后两项逐项抵消：
\begin{equation*}
 \mathcal K_0^{\mu\nu}\kappa_\nu
 =-\kappa^2\kappa^\mu+\kappa^\mu(\kappa^\nu\kappa_\nu)=0.
\end{equation*}
这就是规范变换 $\mathscr C_\mu\mapsto\mathscr C_\mu+\partial_\mu\chi$ 在 Fourier 空间中的切向方向。

协变规范固定项为
\begin{equation*}
 \Lag_{\rm gf}
 =-\frac{\hbar c}{2\xi}(\partial_\mu\mathscr C^\mu)^2.
\end{equation*}
它只修改纵向二次型，因此总核成为\eqnrefs{eq:u1-gauge-fixed-kernel-dp50}。定义横向、纵向投影
\begin{equation*}
 P_T^{\mu\nu}=\eta^{\mu\nu}-\frac{\kappa^\mu\kappa^\nu}{\kappa^2},
 \qquad
 P_L^{\mu\nu}=\frac{\kappa^\mu\kappa^\nu}{\kappa^2},
\end{equation*}
它们满足 $P_TP_L=0$、$P_T^2=P_T$、$P_L^2=P_L$。于是
\begin{equation*}
 \mathcal K_\xi
 =-\kappa^2 P_T-\frac{\kappa^2}{\xi}P_L.
\end{equation*}
因此除去 $\kappa^2=0$ 的物理传播极点以后，横向和纵向方向都已有非零本征值；规范固定正是通过给纯规范方向指定可逆的二次核来定义传播子。
\end{exbox}

规范固定除去路径积分中的重复计数以后，还需要保证不同规范参数给出相同的物理振幅。为此引入一种同时作用于规范场和鬼场的 Grassmann 奇变换；“Grassmann 奇”表示它会改变变量的玻色/费米奇偶性，并遵守带符号的 Leibniz 法则。这个变换称为 Becchi--Rouet--Stora--Tyutin（BRST）变换。其核心代数性质是幂零性 $s^2=0$：连续作用两次返回零。正是这一性质允许把规范冗余写成同调中的“恰当方向”，从而在量子态空间中识别物理等价类。

规范固定项可以写成一次 BRST 变换；为此引入 Nakanishi--Lautrup 辅助场 $B^a$。所谓\emph{辅助场}是指它没有独立传播的动能项，Euler--Lagrange 方程只是代数方程，因此可以在最后直接消去而不增加物理自由度。为与\eqnrefs{eq:generic-infinitesimal-gauge-dp11} 的规范变换约定一致，固定鬼场归一化为 $c^a=-\vartheta^a/g_{\rm G}$；BRST 微分 $s$ 定义为
\begin{align}
 s\mathscr C_\mu^a
 &=(\mathcal D_\mu^{\rm adj})^{ab}c^b,\\
 sc^a
 &=+\frac {g_{\rm G}}2 f^{abc}c^bc^c,\\
 s\bar c^a
 &=B^a,\\
 sB^a&=0.
 \label{eq:brst-transform-general-dp11}
\end{align}
连续作用两次时，真正需要检查的非平凡对象是鬼场和规范场；这不是两个彼此独立的结论，而是同一算符恒等式 $s^2=0$ 在两类非平凡场上的分量检验。分别有
\begin{align}
 s^2c^a&=0,
 \label{eq:brst-ghost-nilpotent-dp68}\\
 s^2\mathscr C_\mu^a&=0.
 \label{eq:brst-gauge-nilpotent-dp68}
\end{align}
第一式由结构常数的 Jacobi 恒等式与三个鬼场的反对称性给出；第二式还需要把 $\partial_\mu(sc)$ 与 $s\mathscr C_\mu$ 中的导数项逐项相消。再结合 $s\bar c^a=B^a$、$sB^a=0$，可汇总为
\begin{equation}
 \boxed{s^2=0}
 \label{eq:brst-nilpotent-dp11}
\end{equation}
对上述全部场成立。幂零性不是记号便利，而是物理态上同调能够定义的必要条件。

规范固定项与鬼场项可以合并成一个 BRST 恰当项：
\begin{equation}
 \boxed{
 \Lag_{\rm gf+gh}
 =\hbar c\,s\left[
 \bar c^a\left(\partial^\mu\mathscr C_\mu^a+\frac\xi2B^a\right)
 \right].}
 \label{eq:brst-exact-gf-dp11}
\end{equation}
展开后为
\begin{equation}
 \frac{\Lag_{\rm gf+gh}}{\hbar c}
 =B^a\partial^\mu\mathscr C_\mu^a
 +\frac\xi2B^aB^a
 -\bar c^a\partial^\mu(\mathcal D_\mu^{\rm adj})^{ab}c^b,
 \label{eq:brst-gf-expanded-dp11}
\end{equation}
消去没有动力学导数的 $B^a$ 就恢复\eqnrefs{eq:general-covariant-gf-dp11,eq:general-ghost-lagrangian-dp11}。因此协变规范固定不是与规范理论并列的另一套动力学，而是同一规范轨道空间的一种坐标选择。

协变量子化为了保持 Lorentz 协变性，会暂时把横向、纵向、类时规范激发和鬼场都放进计算用的 Fock 空间。因此这个空间包含冗余方向，不能直接等同于物理 Hilbert 空间。

BRST 荷 $Q_B$ 在态空间上实现幂零关系 $Q_B^2=0$。物理态取为 BRST 闭合态，而相差一个 BRST 恰当态的两个闭合态代表同一物理状态：
\begin{equation*}
 Q_B|\mathrm{phys}\rangle=0,
 \qquad
 |\mathrm{phys}\rangle\sim|\mathrm{phys}\rangle+Q_B|\chi\rangle.
\end{equation*}
其中 $Q_B|\chi\rangle$ 只改变规范代表而不改变物理矩阵元。于是物理态空间是核与像之商 $\ker Q_B/\operatorname{im}Q_B$；这正是 BRST 上同调的定义。


对生成泛函作无穷小 BRST 换元时，积分测度与规范固定后的作用量保持不变。对单粒子不可约有效作用量 $\Gamma_{\rm eff}$，这一事实写成 Slavnov--Taylor 泛函恒等式
\begin{equation}
 \boxed{\mathcal S(\Gamma_{\rm eff})=0.}
 \label{eq:st-functional-general-dp11}
\end{equation}
对它作不同阶泛函微分可得到各个顶角函数之间的 Slavnov--Taylor 恒等式。阿贝尔 $U(1)$ 中鬼场脱耦，它退化为 Ward--Takahashi 恒等式；非阿贝尔理论中则同时约束物质场、规范场和鬼场顶角，以及波函数与耦合的重整化常数。

量子化还可能暴露另一类更根本的问题。一个连续变换即使保持经典作用量不变，调节后的函数积分测度或量子有效作用量也未必能同时保持同一对称性。若这种破坏不能通过允许的局域反项消去，就说这个经典对称性在量子理论中具有\emph{量子反常}（quantum anomaly）。对整体对称性，反常可以成为真实的量子效应；对必须用来消除非物理自由度的局域规范对称性，未抵消的规范反常会破坏物理态空间的一致性，因此可接受的基本规范理论必须使所有规范反常相互抵消。轴流反常与手征规范反常的差别，正来自这一一般量子测度效应作用在整体对称性或局域规范对称性上的不同物理角色。



Faddeev--Popov 行列式解决“规范轨道重复积分”的测度问题；鬼场把这个泛函行列式局域化；BRST 幂零性进一步保证规范固定后的扩大场空间能够通过上同调投影回物理态空间。三者属于同一规范冗余问题的不同层次，而不是彼此独立的附加规则。




\begin{derivation}[从\eqnrefs{eq:generic-finite-gauge-transform-dp68}到\eqnrefs{eq:generic-infinitesimal-gauge-dp11}]

正文\eqnrefs{eq:generic-finite-gauge-transform-dp68} 给出了联络在有限局域变换下的变化。把 $U$ 展开到规范参数 $\vartheta^a$ 的一阶，并使用 Lie 代数对易关系，就得到沿规范轨道的无穷小切向量~\eqnrefs{eq:generic-infinitesimal-gauge-dp11}。这一步只是在场空间中线性化有限变换。

从有限规范变换
\begin{equation*}
 \mathscr C_\mu'
 =U\mathscr C_\mu U^{-1}
 +\frac{\ii}{g_{\rm G}}(\partial_\mu U)U^{-1}
\end{equation*}
出发。令 $U=I+\ii\vartheta+\order(\vartheta^2)$、$U^{-1}=I-\ii\vartheta+\order(\vartheta^2)$，则
\begin{align*}
 U\mathscr C_\mu U^{-1}
 &=\mathscr C_\mu+\ii[\vartheta,\mathscr C_\mu]+\order(\vartheta^2),\\
 \frac{\ii}{g_{\rm G}}(\partial_\mu U)U^{-1}
 &=-\frac1{g_{\rm G}}\partial_\mu\vartheta+\order(\vartheta^2).
\end{align*}
因此
\begin{equation*}
 \delta\mathscr C_\mu
 =-\frac1{g_{\rm G}}\partial_\mu\vartheta+\ii[\vartheta,\mathscr C_\mu].
\end{equation*}
写成分量并使用 $[T^a,T^b]=\ii f^{abc}T^c$，得到
\begin{equation*}
 \delta\mathscr C_\mu^a
 =-\frac1{g_{\rm G}}\left(\partial_\mu\vartheta^a-g_{\rm G}f^{acb}\mathscr C_\mu^c\vartheta^b\right)
 =-\frac1{g_{\rm G}}(\mathcal D_\mu^{\rm adj})^{ab}\vartheta^b.
\end{equation*}
\end{derivation}

\begin{derivation}[从\eqnrefs{eq:fp-jacobian-definition-dp68}到\eqnrefs{eq:fp-identity-general-dp11}]

规范固定的数学结构与普通多元积分中的换元完全相同，只是“坐标”变成了函数。先把规范场记为 \(\mathscr C_\mu^a(x)\)，无穷小规范参数记为 \(\vartheta^a(x)\)。沿规范轨道作变换
\[
\mathscr C_\mu
\longmapsto
\mathscr C_\mu^\vartheta .
\]
选定规范条件
\[
F^a[\mathscr C](x)=\omega^a(x).
\]
在一条给定规范轨道上，真正需要积分的是参数 \(\vartheta\)；因此定义映射
\[
\mathcal F_{\mathscr C}:
\vartheta
\longmapsto
F[\mathscr C^\vartheta]-\omega .
\]
若在所研究区域内每条规范轨道与切片 \(F=\omega\) 只有一个局部交点，则
\(\mathcal F_{\mathscr C}\) 在交点附近可逆。它的 Fr\'echet 导数为
\[
\left(
\mathcal M_{\rm FP}[\mathscr C]
\right)^{ab}(x,y)
\equiv
\left.
\frac{\delta F^a[\mathscr C^\vartheta](x)}
{\delta\vartheta^b(y)}
\right|_{\vartheta=0},
\]
即正文\eqnrefs{eq:fp-jacobian-definition-dp68}。

先回忆有限维公式。若 \(y=f(x)\) 且根 \(x_*\) 孤立，
\[
\delta^{(n)}[f(x)]
=
\frac{\delta^{(n)}(x-x_*)}
{|\det(\partial f/\partial x)_{x_*}|}.
\]
对 \(x\) 积分得到
\[
1
=
|\det J_f(x_*)|
\int\dd^n x\,\delta^{(n)}[f(x)].
\]
函数空间的局部换元完全平行：
\[
\mathcal D\!\left(F[\mathscr C^\vartheta]-\omega\right)
=
\det\mathcal M_{\rm FP}[\mathscr C^\vartheta]\,
\mathcal D\vartheta .
\]
于是只要切片在局部唯一，
\begin{align*}
\int\mathcal D\vartheta\,
\delta\!\left[F(\mathscr C^\vartheta)-\omega\right]
&=
\left[
\det\mathcal M_{\rm FP}[\mathscr C]
\right]^{-1},
\end{align*}
从而
\[
\boxed{
1
=
\Delta_{\rm FP}[\mathscr C]
\int\mathcal D\vartheta\,
\delta\!\left[
F(\mathscr C^\vartheta)-\omega
\right],
\qquad
\Delta_{\rm FP}[\mathscr C]
\equiv
\det\mathcal M_{\rm FP}[\mathscr C].
}
\]
这就是正文\eqnrefs{eq:fp-identity-general-dp11}。

为了看清这个抽象 Jacobian 的具体内容，取 Yang--Mills 场的无穷小规范变换
\[
\delta_\vartheta \mathscr C_\mu^a
=
(D_\mu\vartheta)^a
=
\partial_\mu\vartheta^a
+g_s f^{acb}\mathscr C_\mu^c\vartheta^b,
\]
并选 Lorenz 型规范条件
\[
F^a[\mathscr C]
=
\partial^\mu\mathscr C_\mu^a.
\]
则
\begin{align*}
\delta_\vartheta F^a(x)
&=
\partial^\mu
(D_\mu\vartheta)^a(x)\\
&=
\int\dd^4 y\,
\left[
\partial^\mu D_\mu^{ab}(x)
\delta^{(4)}(x-y)
\right]
\vartheta^b(y),
\end{align*}
所以
\[
\mathcal M_{\rm FP}^{ab}(x,y)
=
\partial^\mu D_\mu^{ab}(x)
\delta^{(4)}(x-y)
\]
到所选整体号规。Abelian \(U(1)\) 中 \(f^{abc}=0\)，于是
\[
\mathcal M_{\rm FP}
=
\partial^2
\]
与电磁场本身无关；其行列式只给一个可吸收的常数。非 Abelian 理论中 \(D_\mu\) 含规范场，因此
\(\det\mathcal M_{\rm FP}\) 是场依赖的，这正是后续 Faddeev--Popov ghost 与规范场耦合的来源。

若规范切片与同一轨道出现多个交点，上述“一个根对应一个 Jacobian”的局部换元不再全局成立；这就是 Gribov 问题出现的数学位置。Faddeev--Popov 恒等式因此是局部规范轨道坐标公式，而不是无条件的全局几何恒等式。
\end{derivation}

\begin{derivation}[从\eqnrefs{eq:brst-transform-general-dp11}到\eqnrefs{eq:brst-nilpotent-dp11}]

先检验鬼场。由
\begin{equation*}
 sc^a=\frac{g_{\rm G}}2f^{abc}c^bc^c
\end{equation*}
以及分次 Leibniz 法则
\begin{equation*}
 s(c^bc^c)=(sc^b)c^c-c^b(sc^c),
\end{equation*}
得到
\begin{align*}
 s^2c^a
 &=\frac{g_{\rm G}^2}{4}f^{abc}f^{bde}c^dc^ec^c
 -\frac{g_{\rm G}^2}{4}f^{abc}f^{cde}c^bc^dc^e.
\end{align*}
利用鬼场反对易性把三个鬼场统一排序，并对 $(b,d,e)$ 完全反对称化，
\begin{equation*}
 s^2c^a
 =\frac{g_{\rm G}^2}{12}
 \left(
 f^{abm}f^{mde}
 +f^{adm}f^{meb}
 +f^{aem}f^{mbd}
 \right)c^bc^dc^e=0,
\end{equation*}
其中最后一个等号使用 Lie 代数 Jacobi 恒等式。这就是正文\eqnrefs{eq:brst-ghost-nilpotent-dp68}。

再检验规范场。由
\begin{equation*}
 s\mathscr C_\mu^a
 =\partial_\mu c^a-g_{\rm G}f^{abc}\mathscr C_\mu^bc^c
\end{equation*}
再作用一次 $s$，
\begin{equation*}
 s^2\mathscr C_\mu^a
 =\partial_\mu(sc^a)
 -g_{\rm G}f^{abc}(s\mathscr C_\mu^b)c^c
 -g_{\rm G}f^{abc}\mathscr C_\mu^b(sc^c).
\end{equation*}
第一项化为
\begin{align*}
 \partial_\mu(sc^a)
 &=\frac{g_{\rm G}}2f^{ade}
 [ (\partial_\mu c^d)c^e+c^d\partial_\mu c^e]\\
 &=g_{\rm G}f^{ade}(\partial_\mu c^d)c^e,
\end{align*}
恰好抵消第二项中
$-g_{\rm G}f^{abc}(\partial_\mu c^b)c^c$ 的导数部分。剩余项先保持由 Leibniz 法则产生的系数，得到
\begin{align*}
 s^2\mathscr C_\mu^a
 &=g_{\rm G}^2f^{abc}f^{bde}\mathscr C_\mu^d c^e c^c
 -\frac{g_{\rm G}^2}{2}f^{abc}f^{cde}\mathscr C_\mu^b c^dc^e\\
 &=-\frac{g_{\rm G}^2}{2}\mathscr C_\mu^b
 \left(
 f^{abm}f^{mde}
 +f^{adm}f^{meb}
 +f^{aem}f^{mbd}
 \right)c^dc^e\\
 &=0,
\end{align*}
所以得到正文\eqnrefs{eq:brst-gauge-nilpotent-dp68}。此外
\begin{equation*}
 s^2\bar c^a=sB^a=0,
 \qquad
 s^2B^a=0,
\end{equation*}
故正文\eqnrefs{eq:brst-nilpotent-dp11} 对这组场变量全部成立。
\end{derivation}

\begin{derivation}[从\eqnrefs{eq:brst-exact-gf-dp11}到\eqnrefs{eq:brst-gf-expanded-dp11}]

正文\eqnrefs{eq:brst-exact-gf-dp11} 把规范固定与鬼场项写成一个 BRST 恰当量。这里仅展开一次 BRST 微分，检查 $B$ 场、规范固定条件和 Faddeev--Popov 鬼算符的相对符号，从而得到正文\eqnrefs{eq:brst-gf-expanded-dp11}；随后消去代数辅助场 $B^a$ 的一步留在正文中完成。

令
\begin{equation*}
 \Psi_{\rm gf}
 =\bar c^a\left(\partial^\mu\mathscr C_\mu^a+\frac\xi2B^a\right).
\end{equation*}
因为 $\bar c$ 是 Grassmann 奇对象，分次 Leibniz 法则给
\begin{align*}
 s\Psi_{\rm gf}
 ={}&(s\bar c^a)
 \left(\partial\!\cdot\!\mathscr C^a+\frac\xi2B^a\right)
 -\bar c^a\left[
 \partial^\mu(s\mathscr C_\mu^a)+\frac\xi2sB^a
 \right].
\end{align*}
代入正文\eqnrefs{eq:brst-transform-general-dp11} 中的
$s\bar c^a=B^a$、$sB^a=0$ 与
$s\mathscr C_\mu^a=(\mathcal D_\mu^{\rm adj})^{ab}c^b$，得到
\begin{equation*}
 s\Psi_{\rm gf}
 =B^a\partial^\mu\mathscr C_\mu^a
 +\frac\xi2B^aB^a
 -\bar c^a\partial^\mu(\mathcal D_\mu^{\rm adj})^{ab}c^b.
\end{equation*}
乘回正文\eqnrefs{eq:brst-exact-gf-dp11} 中的整体因子 $\hbar c$，就是正文\eqnrefs{eq:brst-gf-expanded-dp11}。
\end{derivation}

